\scanpage{22}
% lines 1–8
\uline{Remarque} Si 2 inv. sur $S$, alors $f$ est un isomorphisme d'algèbres de Lie, \deleted{$\unsure{\rho_E|f}$} permettant d'identifier $E$ et $E'$ (c'est aussi l'identification habituelle d'un espace quadratique et de son dual), et $\rho_E$ est identifié à $\rho_{E'}$ ou à $\widetilde\rho_E$ ou $\widetilde\rho_{E'}\ldots$ ; $f'$ (qui s'identifie à $2\operatorname{id}_E$) n'a alors guère d'intérêt ! (Ces constatations s'appliquent plus général\textsuperscript{t} chaque fois que $\Delta$ est inv.)

% lines 9–14
Si par contre 2 est nul sur $S$, alors $f$ a comme noyau $\mathscr O_S\mathbin{.}\widetilde H$, qui est alors le \uline{centre} de $E=\mathfrak{sl}(2)$. Quant à $f'$, il a comme noyau $\mathscr O_S X\oplus\mathscr O_S Y$, qui est \deleted{aussi} \inserted{alors} l'Algèbre \uline{dérivée} de $E'=\mathfrak{gp}(1)$, \deleted{et \unsure{comme}} \inserted{qui} est aussi l'image de $f$, alors que l'image

% lines 14a
% slanted
\margin{\uline{NB} $\varphi$ alternée, sa connaissance ne permet pas de retrouver $q$ ; $\varphi'$ \uline{n'est pas} alternée (\unsure{i.e.|resp.} $q'$ n'est pas nulle, comme on l'a déjà \unsure{noté}\ldots).}

% lines 15–19
de $f'$ est le noyau $\mathscr O_S\widetilde H$ de $f$. [Par contre le centre de $\mathfrak{gp}(1)$ est nul -- c'est vrai sur tte base $S$, comme on vérifie facilement $\}$ -- tandis que le \unsure{groupe|gpe} dérivé de $\mathfrak{sl}(2).$ est égal au centre $\mathscr O_S\mathbin{.}\widetilde H$] Ici

% lines 20–24
l'existence de $f'$ (cf plus bas un contexte général) peut avoir un intérêt \ldots. Il faudrait par la suite expliciter la structure des formes de $\mathfrak s=\mathfrak{sl}(2)$ et $\mathfrak g=\mathfrak{gp}(1)$ en car. 2, et se convaincre qu'elle est trop triviale pour offrir

% lines foot
% archive sheet number: 21
